\section{Conclusions}\label{sec:conclusions}

Minimizing the expected trap-intensity cost over a Gaussian prior on residual Zernike aberrations, with a deterministic set of sixteen sigma points and a conjugate-gradient descent from phase-fixed WGS, lowers the 95th percentile of the intensity spread of a $10\times10$ array under residual aberrations by a median factor of \ratioMed{} over three seeds (\ratioMin{} to \ratioMax) and \ratioAllMed{} over \nSeedsAll{} seeds (interval \ratioAllCiLo{} to \ratioAllCiHi), at \effRatioMed{} of the efficiency, in a scalar Fourier model with an ideal device. The gain needs the expected cost (the nominal cost with the same solver gives \ablNominalMed), a design that covers the modes of the residual, and a design amplitude and a target level near the ones of the ensemble. It
does not transfer to larger amplitudes or to modes outside the design, and it adds a nominal spread of about one percent. The first-order theory explains the direction of the gain and overstates its size, as the measured ratio is below the first-order prediction of about \thQLinRatio. For random displacements of the traps the acceptance configuration gives a ratio of \ptAccMed{}, and a displacement prior optimized alone for 100 steps gives \poBLong{} at 0.10 bins at \poBLongEff{} of the efficiency; the spread under displacement has a first-order slope term that the design can lower and a curvature floor that the spot width fixes. The next steps are a measurement of the residual of an instrument to give the prior a source, a per-patch prior, a design that includes the second-order response, and a test on hardware.
